%%%%%%%%%%  scan idx 445  |  folio 432  %%%%%%%%%%
\SGAPrintPage{432}%
\(U\) \uline{of} \(x\), \uline{and a sequence of sections} \(f_1,\ldots,f_p,g_1,\ldots,g_r,h_1,\ldots,h_s\) \uline{of} \(\underline{O}_U\) \uline{such that the Koszul complex associated with} \(f_1,\ldots,f_p,g_1,\ldots,g_r\) \uline{is a resolution of} \(\underline{O}_{Y\cap U}\), \uline{the Koszul complex associated with} \(f_1,\ldots,f_p,h_1,\ldots,h_s\) \uline{a resolution of} \(\underline{O}_{Z\cap U}\), \uline{and the Koszul complex associated with} \(f_1,\ldots,f_p,g_1,\ldots,g_r,h_1,\ldots,h_s\) \uline{a resolution of} \(\underline{O}_{T\cap U}\) (\uline{conditions satisfied in particular in the noetherian case when the sequence} \(f_1,\ldots,f_p,g_1,\ldots,g_r,h_1,\ldots,h_s\) \uline{is regular and} \(Y=V(f_1,\ldots,f_p,g_1,\ldots,g_r)\), \(Z=V(f_1,\ldots,f_p,h_1,\ldots,h_s)\)). \uline{Let} \(\underline{I},\underline{J}\) \uline{be the ideals of} \(Y,Z\). \uline{We then have the following isomorphisms}:
\[
\begin{array}{ll}
\text{i)}& \forall i\ge0,\quad
\underline{\Tor}_i^{\underline{O}_X}(\underline{O}_Y,\underline{O}_Z)\xrightarrow{\sim}
\bigwedge^i\underline{\Tor}_1^{\underline{O}_X}(\underline{O}_Y,\underline{O}_Z);\\[0.8ex]
\text{ii)}&
\underline{\Tor}_1^{\underline{O}_X}(\underline{O}_Y,\underline{O}_Z)\simeq
\underline{I}\cap \underline{J}/\underline{I}.\underline{J}\simeq
\Ker\bigl(i^*(\underline{N}_{Y/X})\times j^*(\underline{N}_{Z/X})\to \underline{N}_{T/X}\bigr),
\end{array}
\]
\uline{where we denote by} \(i\) \uline{the immersion of} \(T\) \uline{into} \(Y\) \uline{and by} \(j\) \uline{the immersion of} \(T\) \uline{into} \(Z\).

\uline{In particular}, \uline{for every regular closed immersion} \(Y\to X\), \uline{we have}
\[
\begin{array}{ll}
\text{i')}& \forall i\ge0,\quad
\underline{\Tor}_i^{\underline{O}_X}(\underline{O}_Y,\underline{O}_Y)\xrightarrow{\sim}
\bigwedge^i\underline{\Tor}_1^{\underline{O}_X}(\underline{O}_Y,\underline{O}_Y);\\[0.8ex]
\text{ii')}&
\underline{\Tor}_1^{\underline{O}_X}(\underline{O}_Y,\underline{O}_Y)\simeq \underline{N}_{Y/X}.
\end{array}
\]

i) We saw in 2.3 that \(\bigoplus_{i\ge0}\underline{\Tor}_i^{\underline{O}_X}(\underline{O}_Y,\underline{O}_Z)\) could be endowed with an \(\underline{O}_T\)-Algebra structure; by 2.2, this algebra is alternating, since the immersion of \(Y\) (or \(Z\)) into \(X\) is regular. There therefore exists one and only one algebra homomorphism:
\sgaStableTarget{sga6.e07.eq.2-5-1}\begin{equation}\tag{2.5.1}
\bigwedge \underline{\Tor}_1^{\underline{O}_X}(\underline{O}_Y,\underline{O}_Z)\longrightarrow
\bigoplus_{i\ge0}\underline{\Tor}_i^{\underline{O}_X}(\underline{O}_Y,\underline{O}_Z),
\end{equation}
extending the identity in degree \(1\).

To see that it is an isomorphism, we can check locally. Thus we work in a neighborhood of a point of \(T\) and consider sections \(f_1,\ldots,f_p,g_1,\ldots,g_r,h_1,\ldots,h_s\) of \(\underline{O}_X\) on this neighborhood, having the properties stated in 2.5. Put \(\underline{E}=(\underline{O}_X)^p\), \(\underline{E}'=(\underline{O}_X)^r\), and let \(u:\underline{E}\to\underline{O}_X\) be the homomorphism defined by the \(f_i\), \(v:\underline{E}'\to\underline{O}_X\) the homomorphism defined by the \(g_j\), and \(\underline{K}_\bullet(u)\) and \(\underline{K}_\bullet(v)\) the associated

%%%%%%%%%%  scan idx 446  |  folio 433  %%%%%%%%%%
\SGAPrintPage{433}%
Koszul complexes. The Koszul complex associated with the sequence \(f_1,\ldots,f_p,g_1,\ldots,g_r\) is isomorphic to the tensor-product complex \(\underline{K}_\bullet(u)\otimes_{\underline{O}_X}\underline{K}_\bullet(v)\). Since it is a resolution of \(\underline{O}_Y\) by free \(\underline{O}_X\)-Modules of finite type, the \(\underline{\Tor}_i^{\underline{O}_X}(\underline{O}_Y,\underline{O}_Z)\) are the homology sheaves of the complex \(\underline{K}_\bullet(u)\otimes_{\underline{O}_X}\underline{K}_\bullet(v)\otimes_{\underline{O}_X}\underline{O}_Z\). We can regard it as the tensor product of the complexes \(\underline{K}_\bullet(u)\) and \(\underline{K}_\bullet(v)\otimes_{\underline{O}_X}\underline{O}_Z\), and compute its homology using the spectral sequences associated with a bicomplex (EGA \(0_{\mathrm{III}}\) (11.3.5)), which abut to the homology of the associated simple complex.

Consider the spectral sequence such that \({}'E^2_{pq}=H^I_p\bigl(H^{II}_q(K_{\bullet\bullet})\bigr)\), in the notation of EGA \(0_{\mathrm{III}}\) (11.3.5), with the first index referring to the complex \(\underline{K}_\bullet(u)\) and the second to the complex \(\underline{K}_\bullet(v)\otimes_{\underline{O}_X}\underline{O}_Z\). Now the complex \(\underline{K}_\bullet(v)\otimes_{\underline{O}_X}\underline{O}_Z\) is the Koszul complex over \(\underline{O}_Z\) associated with the sequence of sections \(\overline g_1,\ldots,\overline g_r\), where the \(\overline g_i\) are the images of the \(g_i\) in \(\underline{O}_Z\), and this complex is a resolution of \(\underline{O}_T\) over \(\underline{O}_Z\) since the immersions \(Z\to X\), and \(T\to Z\to X\) are regular (cf. the proof of 1.7 ii). Since \(\bigwedge^p\underline{E}\) is free over \(\underline{O}_X\), the complex \(\bigwedge^p\underline{E}\otimes_{\underline{O}_X}(\underline{K}_\bullet(v)\otimes_{\underline{O}_X}\underline{O}_Z)\) is a resolution of \(\bigwedge^p\underline{E}\otimes_{\underline{O}_X}\underline{O}_T\). Consequently the spectral sequence degenerates, and the \(\underline{\Tor}_i^{\underline{O}_X}(\underline{O}_Y,\underline{O}_Z)\) are the homology sheaves of the complex \(\underline{K}_\bullet(u)\otimes_{\underline{O}_X}\underline{O}_T\). The differential of \(\underline{K}_\bullet(u)\) is expressed in terms of the \(f_i\), which annihilate \(\underline{O}_T\); consequently the differential of \(\underline{K}_\bullet(u)\otimes_{\underline{O}_X}\underline{O}_T\) is zero, and we therefore obtain the isomorphisms:
\sgaStableTarget{sga6.e07.eq.2-5-2}\begin{equation}\tag{2.5.2}
\underline{\Tor}_i^{\underline{O}_X}(\underline{O}_Y,\underline{O}_Z)
\xrightarrow{\sim}
\underline{O}_T\otimes_{\underline{O}_X}\bigwedge^i\underline{E}
\longrightarrow
\bigwedge^i(\underline{O}_T)^p.
\end{equation}

In particular, for \(i=1\) we find:
\[
\underline{\Tor}_1^{\underline{O}_X}(\underline{O}_Y,\underline{O}_Z)\xrightarrow{\sim}(\underline{O}_T)^p,
\]
whence for every \(i\):
\sgaStableTarget{sga6.e07.eq.2-5-3}\begin{equation}\tag{2.5.3}
\underline{\Tor}_i^{\underline{O}_X}(\underline{O}_Y,\underline{O}_Z)
\xrightarrow{\sim}
\bigwedge^i\underline{\Tor}_1^{\underline{O}_X}(\underline{O}_Y,\underline{O}_Z).
\end{equation}

%%%%%%%%%%  scan idx 447  |  folio 434  %%%%%%%%%%
\SGAPrintPage{434}%
Now it follows from 2.2 that the isomorphism:
\[
\bigwedge(\underline{O}_T)^p\xrightarrow{\sim}
\bigoplus_{i\ge0}\underline{\Tor}_i^{\underline{O}_X}(\underline{O}_Y,\underline{O}_Z)
\]
defined by the isomorphisms (2.5.2) is compatible with the algebra structures of the two sides; hence the same is true of the isomorphism:
\[
\bigwedge\underline{\Tor}_1^{\underline{O}_X}(\underline{O}_Y,\underline{O}_Z)\xrightarrow{\sim}
\bigoplus_{i\ge0}\underline{\Tor}_i^{\underline{O}_X}(\underline{O}_Y,\underline{O}_Z)
\]
defined by the isomorphisms (2.5.3). By uniqueness, the latter is therefore the restriction of the homomorphism (2.5.1) that was defined globally. The homomorphism (2.5.1) is thus an isomorphism.

ii) Consider the exact sequence:
\[
0\longrightarrow \underline{I}\longrightarrow \underline{O}_X\longrightarrow \underline{O}_Y\longrightarrow0,
\]
and the resulting exact \(\underline{\Tor}\) sequence:
\[
0\longrightarrow\underline{\Tor}_1^{\underline{O}_X}(\underline{O}_Y,\underline{O}_Z)
\longrightarrow \underline{I}\otimes_{\underline{O}_X}\underline{O}_Z
\longrightarrow \underline{O}_Z.
\]
It follows that \(\underline{\Tor}_1^{\underline{O}_X}(\underline{O}_Y,\underline{O}_Z)\) is the kernel of the homomorphism:
\[
\underline{I}\otimes_{\underline{O}_X}(\underline{O}_X/\underline{J})\longrightarrow \underline{O}_X/J,
\]
that is, \(\underline{I}\cap \underline{J}/\underline{I}.\underline{J}\).

To prove the last relation, we shall show that there is an exact sequence:
\[
\underline{0}\longrightarrow \underline{I}\cap \underline{J}/\underline{I}.\underline{J}
\longrightarrow i^*(\underline{N}_{Y/X})\times j^*(\underline{N}_{Z/X})
\longrightarrow \underline{N}_{T/X}\longrightarrow0.
\]
The homomorphisms \(i^*(\underline{N}_{Y/X})\to \underline{N}_{T/X}\) and \(j^*(\underline{N}_{Z/X})\to \underline{N}_{T/X}\) were defined in (1.6.1); here we take their difference. Furthermore, \(\underline{I}\cap \underline{J}\) maps naturally into \(\underline{I}/\underline{I}^2\), hence into \(i^*(\underline{I}/\underline{I}^2)\), and likewise into \(j^*(\underline{J}/\underline{J}^2)\), and these homomorphisms vanish on \(\underline{I}.\underline{J}\), since \(\underline{I},\underline{J}\subset \underline{K}\) (where \(\underline{K}\) denotes the ideal of \(T\)), defining the arrow on the left.

%%%%%%%%%%  scan idx 448  |  folio 435  %%%%%%%%%%
\SGAPrintPage{435}%
To check exactness of the sequence, we can work locally on \(T\) and use the sections defined in the statement of 2.5. We then have a basis of \(\underline{N}_{T/X}\) formed by the images of the \(f_i,g_j,h_k\), a basis of \(i^*(\underline{N}_{Y/X})\) formed by the images of the \(f_i,g_j\), and a basis of \(j^*(\underline{N}_{Z/X})\) formed by the images of the \(f_i,h_k\). It is then clear that \(i^*(\underline{N}_{Y/X})\times j^*(\underline{N}_{Z/X})\to \underline{N}_{T/X}\) is surjective, and that its kernel consists of the elements whose coordinates on the \((\overline g_j,0)\) and the \((0,\overline h_k)\) vanish, and whose coordinates on \((\overline f_i,0)\) and \((0,\overline f_i)\) are equal for every \(i\). Consequently the sequence is also exact in the middle. The three \(\underline{O}_T\)-Modules involved are locally free of finite type; since \(\underline{\Tor}_1^{\underline{O}_X}(\underline{O}_Y,\underline{O}_Z)\) has rank \(p\) by the proof of i), we see immediately that the sequence is exact.

\sgaStableTarget{sga6.e07.para.2-6}\subsection*{2.6.}
We now suppose that \(X\) is a scheme having an ample sheaf. Every closed subscheme \(Y\) of \(X\) then has an ample sheaf. We know that in this case ``global resolutions exist'', and that the \(K^{\bullet}\) formed respectively with locally free sheaves of finite type and with perfect complexes are identical (II, 2.2.9). If moreover the immersion \(i:Y\to X\) is regular, it is a perfect morphism, and we can define an additive homomorphism \(i_*:K^{\bullet}(Y)\longrightarrow K^{\bullet}(X)\) (IV 2.12) by associating with the class of a perfect complex \(\underline{L}^\bullet\) the class of the complex \(\mathbb Ri_*(\underline{L}^\bullet)=i_*(\underline{L}^\bullet)\) (\(i_*\) being an exact functor).

\textbf{Proposition 2.7.} \uline{Let} \(X\) \uline{be a scheme having an ample sheaf}, \(i:Y\to X\) \uline{a regular closed immersion}, \(\underline{N}\) \uline{the conormal sheaf of} \(Y\) \uline{in} \(X\). \uline{Then}, \uline{for every} \(x\in K^{\bullet}(Y)\), \uline{we have the relation}:
\sgaStableTarget{sga6.e07.eq.2-7-1}\begin{equation}\notag
i^*(i_*(x))=x.\lambda_{-1}(N),
\end{equation}
\uline{where} \(N\) \uline{is the class of} \(\underline{N}\) \uline{in} \(K^{\bullet}(Y)\), \uline{and} \(\lambda_{-1}(N)=\sum_{i\ge0}(-1)^i\lambda^i(N)\).

%%%%%%%%%%  scan idx 449  |  folio 436  %%%%%%%%%%
\SGAPrintPage{436}%
Since both sides of the relation to be proved are additive, we may suppose that \(x\) is the class of an \(\underline{O}_Y\)-Module locally free of finite type \(\underline{M}\). The element \(i^*(i_*(x))\) is then the class of the complex \(\mathbb L i^*(\mathbb Ri_*(\underline{M}))\), obtained by taking a projective resolution of the complex \(\mathbb Ri_*(\underline{M})=i_*(\underline{M})\) and tensoring it with \(\underline{O}_Y\). The homology sheaves of this complex are therefore the \(\underline{\Tor}_i^{\underline{O}_X}(\underline{M},\underline{O}_Y)\). By 2.4, we have the isomorphism:
\[
\underline{\Tor}_i^{\underline{O}_X}(\underline{M},\underline{O}_Y)
\xrightarrow{\sim}
\underline{M}\otimes_{\underline{O}_Y}\underline{\Tor}_i^{\underline{O}_X}(\underline{O}_Y,\underline{O}_Y).
\]
By 2.5, we have the isomorphism:
\[
\underline{\Tor}_i^{\underline{O}_X}(\underline{O}_Y,\underline{O}_Y)
\xrightarrow{\sim}
\bigwedge^i\underline{N}.
\]

The \(\underline{\Tor}_i^{\underline{O}_X}(\underline{M},\underline{O}_Y)\) are therefore \(\underline{O}_Y\)-Modules locally free of finite type, and the class in \(K^{\bullet}(Y)\) of the complex \(\mathbb L i^*(\mathbb Ri_*(\underline{M}))\) is the alternating sum of the classes of its homology sheaves, namely:
\[
i^*(i_*(x))=\sum_{i\ge0}(-1)^i
\operatorname{cl}\bigl(\underline{M}\otimes_{\underline{O}_Y}\underline{\Tor}_i^{\underline{O}_X}(\underline{O}_Y,\underline{O}_Y)\bigr)
=\sum_{i\ge0}(-1)^i x.\lambda^i(N).
\]

\sgaStableTarget{sga6.e07.cor.2-8}\textbf{Corollary 2.8.} \uline{Under the hypotheses of} 2.7, \uline{let} \(x,y\in K^{\bullet}(Y)\); \uline{we have the relation}:
\sgaStableTarget{sga6.e07.eq.2-8-1}\begin{equation}\notag
i_*(xy\lambda_{-1}(N))=i_*(x).i_*(y).
\end{equation}

By 2.7, \(y\lambda_{-1}(N)=i^*(i_*(y))\).

Applying the projection formula, we therefore find:
\[
i_*(xy\lambda_{-1}(N))
=i_*(x.i^*(i_*(y)))
=i_*(x).i_*(y).
\]

\sgaStableTarget{sga6.e07.sec.3}\section*{3. Computation of \texorpdfstring{$K^{\bullet}$}{K} of a blown-up scheme}

\sgaStableTarget{sga6.e07.para.3-1}\textbf{3.1.} In this section, \(i:Y\to X\) denotes a regular closed immersion. Denote by \(X'\) the scheme obtained by blowing up \(Y\), and put \(Y'=Y\times_X X'\).
